\chapter{Types, Saturation, and Ultraproducts}
\label{ch:types-ultraproducts}

\section{Types over parameter sets}
\label{sec:types-ultraproducts-01}

Let \(\mathcal M\) be a \(\Lang\)-structure and \(A\subseteq M\).  Use the
parameter expansion \(\Lang(A)\) and \(\mathcal M_A\) fixed in
\cref{def:parameter-expansion}.

\begin{definition}[Type]
Let \(\bar x=(x_1,\ldots,x_n)\) be a finite tuple of distinct variables.  An
\emph{\(n\)-type over \(A\) consistent with \(\mathcal M\)} is a set
\(p(\bar x)\) of \(\Lang(A)\)-formulas whose free variables are among
\(\bar x\) such that
\[
  \Th_{\Lang(A)}(\mathcal M_A)\cup
  \set{\varphi(\bar c):\varphi(\bar x)\in p}
\]
is satisfiable, where \(\bar c\) is a new tuple of constants.  The type is
\emph{complete} iff, for every \(\Lang(A)\)-formula \(\varphi(\bar x)\),
exactly one of \(\varphi\) and \(\neg\varphi\) belongs to \(p\).  The
phrase \emph{consistent with} is the standard model-theoretic term and is
defined here semantically, by satisfiability; by soundness and completeness
of the fixed first-order calculus, this sense agrees with consistency in
the registered derivability sense.  That departure is declared once here
and governs types over parameter sets and over theories.
\end{definition}

\begin{definition}[Types over a theory and type spaces]
\label{def:type-space}
Let \(T\) be a satisfiable \(\Lang\)-theory.  An \emph{\(n\)-type
consistent with \(T\)} is a set \(p(\bar x)\) of \(\Lang\)-formulas whose
free variables are among the distinct variables in \(\bar x\), such that
\[
  T\cup\{\varphi(\bar c):\varphi(\bar x)\in p\}
\]
is satisfiable for a new \(n\)-tuple of constants \(\bar c\).  It is complete
when it contains exactly one of \(\varphi\) and \(\neg\varphi\) for every
\(\Lang\)-formula \(\varphi(\bar x)\).  Let \(S_n(T)\) be the set of all
complete \(n\)-types consistent with \(T\).  For each formula
\(\varphi(\bar x)\), put
\[
  [\varphi]=\{p\in S_n(T):\varphi\in p\}.
\]
The sets \([\varphi]\) form a basis of clopen sets because
\([\varphi]\cap[\psi]=[\varphi\land\psi]\) and
\(S_n(T)\setminus[\varphi]=[\neg\varphi]\).  The set \(S_n(T)\) with this
topology is the \emph{\(n\)-type space of \(T\)}.  The topology is required
to express isolation and compactness of complete descriptions; it is not the
set of tuples realized in one chosen model.  This distinction is used below
in the omitting-types theorem and the Ryll--Nardzewski theorem.
\end{definition}

\begin{theorem}[Compactness of type spaces]
\label{thm:type-space-compact}
Assume the Axiom of Choice in the ambient set theory.  If \(T\) is satisfiable, then
\(S_n(T)\) is compact, Hausdorff, and has a
basis of clopen sets.
\textit{Status: complete proof.}
\end{theorem}

\begin{proof}
The displayed basis in \cref{def:type-space} consists of clopen sets.  If
\(p\ne q\), choose \(\varphi\in p\setminus q\); completeness gives
\(\neg\varphi\in q\), so the disjoint neighborhoods \([\varphi]\) and
\([\neg\varphi]\) separate them.

For compactness it suffices to consider a cover by basic open sets: an
arbitrary open cover refines to one by basic sets, because every point of an
open set lies in a basic set contained in that open set, and a finite
subfamily of the basic refinement yields a finite subcover of the original
cover upon replacing each basic set by a member of the original cover
containing it.

So let \(\set{[\varphi_i]:i\in I}\) cover \(S_n(T)\) and suppose, toward a
contradiction, that no finite subfamily covers \(S_n(T)\).  Let \(\bar c\)
be a new \(n\)-tuple of constants and consider the set
\[
  T\cup\set{\neg\varphi_i(\bar c):i\in I}.
\]
Every finite subset of the displayed set is contained in
\(T\cup\set{\neg\varphi_j(\bar c):j\in J}\) for some finite \(J\subseteq I\).
Since \(\set{[\varphi_j]:j\in J}\) is a finite subfamily, it does not cover
\(S_n(T)\), so there is a complete type \(p\in S_n(T)\) with
\(\varphi_j\notin p\) for every \(j\in J\); completeness of \(p\) gives
\(\neg\varphi_j\in p\).  By \cref{def:type-space},
\(T\cup\set{\psi(\bar c):\psi\in p}\) is satisfiable, and any model of it is
a model of \(T\cup\set{\neg\varphi_j(\bar c):j\in J}\).

By first-order compactness (\cref{thm:fol-compactness}, clause 1, applied in
the signature expanded by the constants \(\bar c\)), the displayed set has a
model \(\mathcal N\); let \(\bar b\) be the interpretation of \(\bar c\).
The complete-type step, which here plays the role of the Lindenbaum lemma,
extracts a point of \(S_n(T)\) from \(\mathcal N\): let \(q\) be the set of
all \(\Lang\)-formulas \(\psi(\bar x)\) with free variables among \(\bar x\)
such that \(\mathcal N\entails\psi(\bar b)\).  For every \(\Lang\)-formula
\(\psi(\bar x)\), exactly one of \(\psi\) and \(\neg\psi\) is satisfied by
\(\bar b\) in \(\mathcal N\), so \(q\) contains exactly one of them; and
\(\mathcal N\) itself is a model of \(T\cup\set{\psi(\bar c):\psi\in q}\),
because it satisfies \(T\) and interprets \(\bar c\) as \(\bar b\).  Hence
\(q\) is a complete \(n\)-type consistent with \(T\), that is,
\(q\in S_n(T)\).  For every \(i\in I\), the model \(\mathcal N\) satisfies
\(\neg\varphi_i(\bar b)\), so \(\varphi_i\notin q\) and
\(q\notin[\varphi_i]\).  This contradicts the assumption that
\(\set{[\varphi_i]:i\in I}\) covers \(S_n(T)\).
\end{proof}

\begin{example}[A finite type space]
Let \(T=\Th(\mathcal M)\), where \(\mathcal M\) is a two-element structure in
the equality-only signature.  Every model of \(T\) has exactly two
elements and is isomorphic to \(\mathcal M\), so every complete type
consistent with \(T\) is realized in \(\mathcal M\).  Its one-type space
has one point, since the two
elements are exchanged by an automorphism.  Its two-type space has exactly two
points: the complete type containing \(x_1=x_2\) and the complete type
containing \(x_1\ne x_2\).  Both points are clopen, so this type space is
discrete.
\end{example}

